\chapter{\ifenglish Introduction\else บทนำ\fi}

\section{\ifenglish Project rationale\else ที่มาของโครงงาน\fi}

ในปัจจุบัน การจัดการสินค้าประเภทสกรูในร้านค้าและคลังสินค้าจำเป็นต้องอาศัยการตรวจสอบ จำแนกประเภท และตรวจนับจำนวนสินค้า เพื่อให้สามารถบริหารจัดการสินค้าคงคลังได้อย่างถูกต้อง อย่างไรก็ตาม สกรูบางประเภทมีรูปร่างและขนาดใกล้เคียงกัน ทำให้การตรวจสอบด้วยสายตาและการนับด้วยมืออาจเกิดความผิดพลาดได้ โดยเฉพาะเมื่อต้องจัดการกับสินค้าจำนวนมาก ซึ่งอาจส่งผลให้การดำเนินงานใช้เวลานานและเกิดความคลาดเคลื่อนในการบันทึกข้อมูลสินค้า

จากปัญหาดังกล่าว โครงงานนี้จึงมีแนวคิดในการพัฒนาระบบตรวจจับและวิเคราะห์สกรูหัวหกเหลี่ยมอัตโนมัติ โดยประยุกต์ใช้เทคนิคการประมวลผลภาพ (Image Processing) ร่วมกับกล้อง Webcam และกล่องควบคุมแสง (Controlled Lighting Box) เพื่อช่วยลดผลกระทบจากแสงภายนอกและให้ได้ภาพที่มีสภาพแวดล้อมสม่ำเสมอ จากนั้นนำภาพเข้าสู่กระบวนการเตรียมภาพ (Image Preprocessing) การแบ่งส่วนภาพ (Segmentation) และการติดตามเส้นขอบวัตถุ (Edge Following) เพื่อวิเคราะห์ลักษณะและขอบเขตของสกรู รวมถึงสนับสนุนการตรวจนับและวิเคราะห์ขนาดของวัตถุ

นอกจากนี้ โครงงานยังมีแนวทางในการพัฒนาเว็บแอปพลิเคชัน (Web Application) สำหรับแสดงผลการวิเคราะห์และจัดการข้อมูลสินค้า เพื่อสนับสนุนการทำงานของพนักงานขาย พนักงานคลังสินค้า และเจ้าของร้าน โดยมุ่งหวังให้ระบบช่วยลดเวลาในการตรวจสอบ ลดความผิดพลาดจากการทำงานด้วยมือ และเพิ่มประสิทธิภาพในการตรวจนับและบริหารจัดการสินค้าคงคลัง

\section{\ifenglish Objectives\else วัตถุประสงค์ของโครงงาน\fi}

\begin{enumerate}                            
    \item เพื่อพัฒนาระบบตรวจจับและจำแนกประเภทสกรูหัวหกเหลี่ยมโดยใช้เทคนิคการประมวลผลภาพ
    \item เพื่อศึกษาและพัฒนากระบวนการวิเคราะห์รูปร่างและขนาดของสกรูจากภาพถ่าย
    \item เพื่อพัฒนาระบบตรวจนับจำนวนสกรูจากภาพภายใต้เงื่อนไขการจัดวางวัตถุที่กำหนด
    \item เพื่อออกแบบและพัฒนาเว็บแอปพลิเคชันสำหรับแสดงผลและจัดการข้อมูลสินค้า
    \item เพื่อศึกษาความเป็นไปได้ในการนำระบบประมวลผลภาพมาประยุกต์ใช้กับการจัดการสินค้าคงคลัง
\end{enumerate}

\section{\ifenglish Project scope\else ขอบเขตของโครงงาน\fi}

\subsection{\ifenglish Data scope\else ขอบเขตด้านข้อมูล\fi}

โครงงานนี้มุ่งเน้นการตรวจจับและจำแนกสกรูหัวหกเหลี่ยมจำนวน 10 Classes ตามประเภทที่กำหนดไว้ โดยใช้ภาพถ่ายของสกรูที่จัดเตรียมขึ้นสำหรับการพัฒนาระบบ

การเก็บข้อมูลภาพจะดำเนินการภายใต้สภาพแวดล้อมที่ควบคุมได้ โดยใช้กล่องควบคุมแสงเพื่อช่วยลดผลกระทบจากแสงภายนอก และกำหนดตำแหน่งการจัดวางสกรูให้เหมาะสมต่อการถ่ายภาพ

ภาพที่นำมาใช้ในการพัฒนาระบบจะต้องมีลักษณะและคุณภาพที่เหมาะสมต่อการประมวลผล เช่น มองเห็นขอบเขตของสกรูได้ชัดเจน ไม่เกิดการซ้อนทับกัน และมีความแตกต่างระหว่างวัตถุกับพื้นหลังเพียงพอสำหรับการแยกวัตถุ

เนื่องจากโครงงานมีข้อจำกัดด้านระยะเวลาและทรัพยากรในการดำเนินงาน จึงกำหนดจำนวนประเภทของสกรูไว้ที่ 10 Classes โดยมุ่งเน้นการเก็บข้อมูลให้ครอบคลุมลักษณะสำคัญของสกรูแต่ละประเภท มากกว่าการเพิ่มจำนวนประเภทของข้อมูลโดยไม่มีขอบเขตที่ชัดเจน

\subsection{\ifenglish Hardware scope\else ขอบเขตด้านฮาร์ดแวร์\fi}

\begin{enumerate}                            
    \item ออกแบบและจัดทำกล่องควบคุมแสง (Controlled Lighting Box) สำหรับใช้ในการถ่ายภาพสกรู โดยมีขนาด 40 × 50 × 40 เซนติเมตร เพื่อช่วยควบคุมสภาพแวดล้อมในการถ่ายภาพและลดผลกระทบจากแสงภายนอก
    \item กำหนดพื้นที่สำหรับจัดวางชิ้นงานภายในกล่อง โดยใช้พื้นที่ขนาด A3 (29.7 × 42 เซนติเมตร) เพื่อรองรับการจัดวางสกรูสำหรับการถ่ายภาพและประมวลผล
    \item ใช้กล้อง Webcam Logitech C920e ความละเอียด 1920 × 1080 พิกเซล สำหรับรับภาพสกรูเข้าสู่ระบบ โดยติดตั้งกล้องในตำแหน่งที่เหมาะสมกับพื้นที่ถ่ายภาพ เพื่อให้สามารถรับภาพชิ้นงานได้อย่างชัดเจน
    \item กำหนดรูปแบบการจัดวางสกรูให้อยู่ภายในพื้นที่ถ่ายภาพ โดยแยกชิ้นงานออกจากกันและหลีกเลี่ยงการวางซ้อน เพื่อให้เหมาะสมกับกระบวนการตรวจจับและวิเคราะห์ภาพ
\end{enumerate}

\subsection{\ifenglish Technical scope\else ขอบเขตด้านเทคนิคและอัลกอริทึม\fi}

โครงงานนี้มุ่งเน้นการประยุกต์ใช้ Image Processing สำหรับตรวจจับและวิเคราะห์ลักษณะของสกรู สำหรับการวิเคราะห์ขนาดในหน่วยจริง จะมีการพิจารณาการใช้วัตถุอ้างอิง (Reference) ที่ทราบขนาด เพื่อใช้ในการปรับเทียบระหว่างหน่วยพิกเซลและหน่วยความยาวจริง ทั้งนี้ ความถูกต้องของการวัดจะขึ้นอยู่กับระยะห่าง มุมกล้อง และสภาพแวดล้อมในการถ่ายภาพ

\subsection{\ifenglish System and usage scope\else ขอบเขตด้านระบบและการใช้งาน\fi}

ระบบที่พัฒนาขึ้นมีองค์ประกอบหลัก ได้แก่ ส่วนรับภาพ ส่วนประมวลผลภาพ และส่วนแสดงผลผ่านเว็บแอปพลิเคชัน

ระบบมีแนวทางรองรับการรับภาพจาก Webcam เพื่อใช้ในการประมวลผล โดยมีการออกแบบกล่องควบคุมแสงเพื่อช่วยสร้างสภาพแวดล้อมในการถ่ายภาพที่เหมาะสม

ในส่วนของเว็บแอปพลิเคชัน มีขอบเขตการใช้งานดังนี้

\begin{enumerate}                            
    \item รองรับการรับภาพจากกล้องหรือการนำเข้าภาพตามรูปแบบที่ระบบกำหนด
    \item แสดงภาพและผลลัพธ์จากกระบวนการประมวลผล
    \item แสดงข้อมูลที่เกี่ยวข้องกับประเภทและจำนวนของสกรู
    \item สนับสนุนการจัดการข้อมูลสินค้าและข้อมูลคลังสินค้า
    \item ออกแบบส่วนติดต่อผู้ใช้ให้เหมาะสมกับพนักงานขาย พนักงานคลังสินค้า และเจ้าของร้าน
\end{enumerate}

การทำงานแบบเรียลไทม์ผ่าน Webcam เป็นแนวทางการพัฒนาของระบบ โดยประสิทธิภาพและความเร็วในการประมวลผลจะต้องได้รับการทดสอบเพิ่มเติมตามความสามารถของฮาร์ดแวร์และซอฟต์แวร์ที่เลือกใช้

\subsection{ขอบเขตด้านความแม่นยำและข้อจำกัด}

ประสิทธิภาพของระบบประมวลผลภาพอาจได้รับผลกระทบจากปัจจัยหลายประการ เช่น คุณภาพของภาพ ความสว่างของแสง การสะท้อนของพื้นผิวโลหะ ความแตกต่างระหว่างสกรูและพื้นหลัง รวมถึงการซ้อนทับกันของวัตถุ

ในกรณีที่สกรูมีลักษณะใกล้เคียงกัน ระบบอาจจำแนกประเภทได้ยาก โดยเฉพาะเมื่อมีความแตกต่างด้านขนาดเพียงเล็กน้อย นอกจากนี้ หากสกรูซ้อนทับกัน อาจทำให้เส้นขอบของวัตถุเชื่อมติดกันและส่งผลต่อการตรวจนับจำนวน

โครงงานจึงกำหนดแนวทางการจัดวางสกรูให้แยกจากกันอย่างชัดเจน และควบคุมสภาพแวดล้อมในการถ่ายภาพให้เหมาะสม เพื่อช่วยลดปัจจัยรบกวนในการประมวลผล

สำหรับการวิเคราะห์ขนาด ความถูกต้องของผลลัพธ์ขึ้นอยู่กับการปรับเทียบภาพและความคงที่ของระยะห่างระหว่างกล้องกับวัตถุ หากเงื่อนไขการถ่ายภาพเปลี่ยนแปลง อาจจำเป็นต้องปรับเทียบระบบใหม่

ทั้งนี้ ประสิทธิภาพของระบบจะต้องได้รับการประเมินจากการทดลองจริง โดยพิจารณาความถูกต้องในการจำแนกประเภท การตรวจนับ และการวิเคราะห์ขนาด ก่อนสรุปผลการทำงานของระบบ

\section{\ifenglish Expected outcomes\else ประโยชน์ที่ได้รับ\fi}

\begin{enumerate}                            
    \item ช่วยลดเวลาในการตรวจสอบและตรวจนับสกรูเมื่อเทียบกับการดำเนินงานด้วยมือ
    \item ช่วยลดภาระงานของพนักงานขายและพนักงานคลังสินค้าในการตรวจสอบและจำแนกประเภทสินค้า
    \item ช่วยลดโอกาสเกิดความผิดพลาดจากการจำแนกประเภทและบันทึกข้อมูลด้วยมือ
    \item ช่วยสนับสนุนการจัดเก็บและจัดการข้อมูลสินค้าให้เป็นระบบมากขึ้น
    \item ช่วยให้เจ้าของร้านสามารถนำข้อมูลสินค้าไปใช้ประกอบการวางแผนและบริหารจัดการสินค้าคงคลัง
    \item เป็นแนวทางในการประยุกต์ใช้เทคนิคการประมวลผลภาพกับงานตรวจสอบและจัดการสินค้าประเภทอื่นในอนาคต
\end{enumerate}

\section{\ifenglish Technology and tools\else เทคโนโลยีและเครื่องมือที่ใช้\fi}

\subsection{\ifenglish Hardware technology\else เทคโนโลยีด้านฮาร์ดแวร์\fi}

\begin{enumerate}                            
    \item \textbf{Webcam:} ใช้สำหรับรับภาพสกรูเข้าสู่ระบบประมวลผล โดยกล้องที่พิจารณาใช้ในโครงงานคือ Logitech C920e ซึ่งมีความละเอียด 1920 × 1080 พิกเซล
    \item \textbf{Controlled Lighting Box:} กล่องควบคุมแสงสำหรับลดผลกระทบจากแสงภายนอกและช่วยสร้างสภาพแวดล้อมในการถ่ายภาพที่สม่ำเสมอ
    \item \textbf{Lighting System:} ระบบให้แสงสว่างภายในกล่อง เพื่อช่วยให้ภาพมีความสว่างเหมาะสมต่อการประมวลผล
    \item \textbf{Reference Object:} วัตถุอ้างอิง เช่น ไม้บรรทัด สำหรับใช้ในการปรับเทียบขนาดจากภาพในกรณีที่ต้องการวิเคราะห์ขนาดจริง
\end{enumerate}

\subsection{\ifenglish Software technology\else เทคโนโลยีด้านซอฟต์แวร์\fi}

ซอฟต์แวร์ที่ใช้ในการพัฒนาระบบประกอบด้วยเครื่องมือสำหรับประมวลผลภาพ พัฒนาเว็บแอปพลิเคชัน และจัดการข้อมูล โดยจะเลือกใช้ให้เหมาะสมกับข้อกำหนดของโครงงาน

\begin{enumerate}                            
    \item \textbf{Frontend:} พัฒนาด้วย React ร่วมกับ TypeScript และ Vite ทำหน้าที่แสดงหน้าจอการใช้งาน รับคำสั่งจากผู้ใช้งาน และแสดงผลการตรวจสอบ
    \item \textbf{Backend:} พัฒนาด้วย Node.js ทำหน้าที่รับส่งข้อมูลระหว่าง Frontend กระบวนการประมวลผลภาพ และฐานข้อมูล
    \item \textbf{Database:} ใช้ PostgreSQL สำหรับจัดเก็บข้อมูลสินค้าคงคลัง ประวัติการตรวจสอบ และข้อมูลที่เกี่ยวข้องกับการดำเนินงานของระบบ
    \item \textbf{OpenCV:} ใช้สำหรับการรับภาพจาก Webcam และการประมวลผลภาพ
\end{enumerate}

\section{\ifenglish Project plan\else แผนการดำเนินงาน\fi}

\begin{plan}{8}{2026}{3}{2027}
    
    % 1. ศึกษาวิเคราะห์ข้อมูลและความต้องการของระบบ
    \planitem{8}{2026}{8}{2026}{ศึกษาวิเคราะห์ข้อมูลและความต้องการของระบบ}
    
    % 2. จัดเตรียมและสร้าง Dataset
    \planitem{9}{2026}{10}{2026}{จัดเตรียมและสร้าง Dataset}
    
    % 3. พัฒนาเว็บแอปพลิเคชัน
    \planitem{10}{2026}{12}{2026}{พัฒนาเว็บแอปพลิเคชัน}
    
    % 4. พัฒนาระบบคลังสินค้าและฐานข้อมูล
    \planitem{11}{2026}{1}{2027}{พัฒนาระบบคลังสินค้าและฐานข้อมูล}
    
    % 5. พัฒนาและฝึกสอนโมเดล AI
    \planitem{11}{2026}{2}{2027}{พัฒนาและฝึกสอนโมเดล AI}
    
    % 6. ทดลองและปรับโมเดลให้แม่นยำมากขึ้น
    \planitem{11}{2026}{2}{2027}{ทดลองและปรับโมเดลให้แม่นยำมากขึ้น}
    
    % 7. ทดสอบและเชื่อมต่อระบบ
    \planitem{1}{2027}{2}{2027}{ทดสอบและเชื่อมต่อระบบ}
    
    % 8. ประเมินผลและปรับปรุงระบบ
    \planitem{2}{2027}{3}{2027}{ประเมินผลและปรับปรุงระบบ}
    
    % 9. จัดทำเอกสารและนำเสนอ
    \planitem{3}{2027}{3}{2027}{จัดทำเอกสารและนำเสนอ}
    
\end{plan}

\section{\ifenglish Roles and responsibilities\else บทบาทและความรับผิดชอบ\fi}

สมาชิกในโครงงานแบ่งหน้าที่ความรับผิดชอบตามความถนัดของแต่ละบุคคล โดยมีการทำงานร่วมกันและช่วยสนับสนุนงานในส่วนต่าง ๆ เพื่อให้การดำเนินโครงงานเป็นไปอย่างมีประสิทธิภาพ มีรายละเอียดดังนี้

\textbf{1. นางสาวณัฐรดา หนูจิตร รหัส 660610757}

รับผิดชอบหลักด้านการออกแบบและจัดทำฮาร์ดแวร์ของโครงงาน รวมถึงการจัดเตรียมและจัดทำชุดข้อมูล (Dataset) โดยมีส่วนร่วมในการทำงานทั้งด้านฮาร์ดแวร์และซอฟต์แวร์ นอกจากนี้ยังช่วยประสานงานนัดหมายและสนับสนุนการดำเนินงานด้านต่าง ๆ ของโครงงาน

\textbf{2. นางสาวทิพวรินท์ สีห์วรางกูร รหัส 660610760}

รับผิดชอบหลักด้านการออกแบบและพัฒนาเว็บไซต์ของระบบ รวมถึงการจัดทำฟังก์ชันและส่วนติดต่อผู้ใช้ โดยร่วมมือกับสมาชิกคนอื่น ๆ ในการดำเนินงานส่วนต่าง ๆ ของโครงงาน

\textbf{3. นายสิรวิชญ์ สีทอง รหัส 660610802}

รับผิดชอบหลักด้านการพัฒนาและดำเนินงานเกี่ยวกับการประมวลผลภาพที่เกี่ยวข้อง โดยร่วมมือกับสมาชิกคนอื่น ๆ ในการดำเนินงานส่วนต่าง ๆ ของโครงงาน

ทั้งนี้ สมาชิกทุกคนมีส่วนร่วมในการจัดทำรายงานและสไลด์นำเสนอ รวมถึงช่วยเหลือและสนับสนุนงานของสมาชิกคนอื่น ๆ ตามความเหมาะสม เพื่อให้โครงงานสำเร็จตามวัตถุประสงค์ที่กำหนด

\section{\ifenglish Social, health, safety, legal, and cultural impacts\else ผลกระทบด้านสังคม สุขภาพ ความปลอดภัย กฎหมาย และวัฒนธรรม\fi}

\subsection{\ifenglish Social impacts\else ผลกระทบด้านสังคม\fi}

โครงงานมีแนวทางในการช่วยลดภาระงานของพนักงานขายและพนักงานคลังสินค้า โดยนำเทคโนโลยีการประมวลผลภาพมาช่วยตรวจสอบและจัดการข้อมูลสินค้า ซึ่งอาจช่วยให้การทำงานมีความสะดวกและเป็นระบบมากขึ้น

\subsection{\ifenglish Health impacts\else ผลกระทบด้านสุขภาพ\fi}

ระบบมีแนวทางช่วยลดการทำงานซ้ำ ๆ จากการตรวจนับและจำแนกสกรูด้วยมือ ซึ่งอาจช่วยลดความเมื่อยล้าจากการปฏิบัติงาน อย่างไรก็ตาม ผลกระทบดังกล่าวยังต้องประเมินจากการทดลองใช้งานจริง

\subsection{\ifenglish Safety impacts\else ผลกระทบด้านความปลอดภัย\fi}

การออกแบบกล่องควบคุมแสงและติดตั้งอุปกรณ์ต้องคำนึงถึงความปลอดภัยในการใช้งาน เช่น ความมั่นคงของโครงสร้างกล่อง การติดตั้งกล้องและระบบไฟ รวมถึงการจัดการสายไฟให้เหมาะสม เพื่อป้องกันอุบัติเหตุระหว่างการใช้งาน

\subsection{\ifenglish Legal impacts\else ผลกระทบด้านกฎหมาย\fi}

การจัดเก็บข้อมูลสินค้าและข้อมูลผู้ใช้งานผ่านเว็บแอปพลิเคชันควรคำนึงถึงความปลอดภัยของข้อมูลและสิทธิ์ในการเข้าถึงข้อมูล โดยเฉพาะหากระบบมีการจัดเก็บข้อมูลส่วนบุคคลของผู้ใช้งาน

\subsection{\ifenglish Cultural impacts\else ผลกระทบด้านวัฒนธรรม\fi}

โครงงานมีแนวทางในการส่งเสริมการนำเทคโนโลยีมาใช้ในการทำงานและการจัดการสินค้า โดยมุ่งเน้นการช่วยสนับสนุนผู้ปฏิบัติงานให้สามารถทำงานได้สะดวกขึ้น ไม่ใช่การทดแทนบทบาทของผู้ใช้งานทั้งหมด แต่เป็นการนำเทคโนโลยีมาช่วยเพิ่มประสิทธิภาพในการทำงาน